\begingroup
\setlength\LTleft{0\linewidth}
\setlength\LTright{0\linewidth}
\fontsize{11.0pt}{13.2pt}\selectfont

\setlength{\LTpost}{0mm}
\begin{longtable}{@{\extracolsep{\fill}}lrrrrrrr}

\toprule
 &  & \multicolumn{3}{c}{\textbf{Ponderação igual}} & \multicolumn{3}{c}{\textbf{Ponderado por preço} (4)} \\ 
\cmidrule(lr){3-5} \cmidrule(lr){6-8}
\textbf{Fase} & \textbf{Dias} & \textbf{Coef.} & \textbf{t} & \textbf{p} & \textbf{Coef.} & \textbf{t} & \textbf{p} \\ 
\midrule\addlinespace[2.5pt]
Dias fora de promoção (1) & 692 & 0,99\% & 16,18 & 0,000 & 0,27\% & 6,17 & 0,000 \\
Pré-promoção (2) & 57 & -0,13\% & -0,76 & 0,448 & -0,24\% & -1,41 & 0,157 \\
Promoção, primeira metade & 64 & -0,59\% & -2,91 & 0,004 & -0,39\% & -2,55 & 0,011 \\
Promoção, segunda metade & 57 & -0,11\% & -0,44 & 0,660 & -0,22\% & -1,72 & 0,086 \\
Pós-promoção (3) & 113 & 0,37\% & 2,11 & 0,035 & 0,28\% & 1,67 & 0,095 \\
Virada do ano & 42 & 0,54\% & 2,25 & 0,025 & 0,60\% & 3,52 & 0,000 \\
\bottomrule
\end{longtable}
\begin{minipage}{\linewidth}
Fonte: Elaboração própria.\\
Regressão do retorno diário de mercado nas variáveis de fase e de virada do ano. Erros-padrão de Newey-West, com defasagens pela regra apresentada por Wooldridge.\\
(1) A constante é a média dos dias fora de promoção.\\
(2) Período anterior à promoção, com metade de sua duração.\\
(3) Período posterior à promoção, com duração igual à da promoção.\\
(4) A ponderação usa o preço da observação anterior.\\
\end{minipage}
\endgroup
